\documentclass{article}

\usepackage[T1]{fontenc}

\usepackage{booktabs,graphicx,siunitx}

\begin{document}

\section{Kendall's tau}

\begin{tabular}{lllll}
\toprule
Pair & $\tau$b [95\% CI] & z & p & N \\
\midrule
motivation\_1 -- motivation\_2 & 0.53 [0.47, 0.60] & 10.64 & $<$.001 & 240 \\
\bottomrule
\end{tabular}

\includegraphics{figures/01_kendalls-tau/figure_scatter.png}

$\tau$ (Kendall’s tau-b) is a rank correlation based on concordant vs. discordant pairs, with a correction for ties - well suited to small samples and many ties. Sign = direction, magnitude = strength; p tests significance. Your significance threshold ($\alpha$) is 0.05.

A Kendall's tau-b correlation gave $\tau$=.53 [.47, .60], p $<$ .001, N=240.

Statistical basis: Kendall's tau-b rank correlation (tie-corrected) (Kendall, M. G. (1938). "A new measure of rank correlation." Biometrika, 30(1/2), 81-93.).

\end{document}